\documentclass[12pt]{article}
\usepackage[top=2.8cm,bottom=2.8cm,left=2.5cm,right=2.5cm]{geometry}
\usepackage[utf8]{inputenc}
\usepackage[T1]{fontenc}
\usepackage{amsmath,amssymb}
\usepackage{booktabs}
\usepackage{graphicx}
\usepackage{enumitem}
\usepackage{setspace}
\usepackage{parskip}
\usepackage{tikz}
\usepackage{pgfplots}
\pgfplotsset{compat=1.18}
\usepackage{fancyhdr}
\setlength{\headheight}{14pt}
\usepackage{xcolor}
\usepackage{mdframed}

\pagestyle{fancy}
\fancyhf{}
\fancyhead[L]{\small ECON 311 --- Fall 2026}

\fancyhead[R]{\small Practice Problem Set 1}
\fancyfoot[C]{\small \thepage}
\renewcommand{\headrulewidth}{0.4pt}

\setlength{\parskip}{6pt}
\onehalfspacing

\begin{document}

%-----------------------------------------------------------------------
% Title block
%-----------------------------------------------------------------------
\begin{center}
  {\LARGE \textbf{Practice Problem Set 1}}\\[10pt]
  {\large \textbf{ECON 311 --- Intermediate Macroeconomics --- Fall 2026}}\\[6pt]

  
\end{center}

\bigskip


\bigskip\hrule\bigskip

%-----------------------------------------------------------------------
% PART I: FOUNDATIONAL PROBLEMS
%-----------------------------------------------------------------------


%-----------------------------------------------------------------------
\noindent\textbf{Problem 1} \quad \textit{Growth-rate arithmetic.}

\medskip
Suppose $x$, $y$, and $w$ grow at constant rates $g_x$, $g_y$, and $g_w$,
respectively.  Express the growth rate of $z$ in terms of these growth rates
for each of the following:

\begin{enumerate}[label=(\alph*),leftmargin=2em,itemsep=4pt]
  \item $\displaystyle z = x^{2} y^{3}$
  \item $\displaystyle z = \frac{x^{0.6}}{y^{0.4}}$
  \item $\displaystyle z = x^{-0.5}\left(\frac{1}{w}\right)^{0.3} y^{0.8}$
  \item $\displaystyle z = \left(\frac{x\,y}{w^{2}}\right)^{1/2}$
\end{enumerate}

\bigskip

%-----------------------------------------------------------------------
\noindent\textbf{Problem 2} \quad \textit{Returns to scale.}

\medskip
For each production function below, determine whether it exhibits
\emph{increasing}, \emph{constant}, or \emph{decreasing} returns to scale
in capital $K$ and labour $L$.  (Treat $\bar{A}$ as a fixed positive
constant throughout.)  Show your work clearly.

\begin{enumerate}[label=(\alph*),leftmargin=2em,itemsep=4pt]
  \item $Y = \bar{A}\,K^{1/4}L^{3/4}$
  \item $Y = K^{0.6}L^{0.6}$
  \item $Y = K^{1/3}L^{2/3} + \bar{A}$
  \item $Y = K + K^{0.5}L^{0.5}$
\end{enumerate}

\bigskip
\newpage
%-----------------------------------------------------------------------
\noindent\textbf{Problem 3} \quad \textit{National income accounting --- three methods.}

\medskip
The economy of \textbf{Silverdale} can be described by the following
activities during one year:

\begin{itemize}[leftmargin=1.5em,itemsep=3pt]
  \item A \textbf{lumber company} hires workers for \$18{,}000 in wages,
        harvests 500 cubic metres of timber, and sells 400 cubic metres to
        a furniture manufacturer at \$60 per cubic metre.  The remaining
        100 cubic metres are milled into finished planks and sold directly
        to households for \$9{,}000.

  \item A \textbf{furniture manufacturer} purchases 400 cubic metres of
        timber at \$60 per cubic metre, pays its workers \$30{,}000 in wages,
        and produces 50 chairs.  Of those chairs: 35 are sold to households
        at \$1{,}200 each, 10 are sold to a government-run community centre
        at \$1{,}200 each, and 5 are exported at \$1{,}400 each.
\end{itemize}

\begin{enumerate}[label=(\alph*),leftmargin=2em,itemsep=4pt]
  \item Calculate GDP using the \textbf{expenditure approach}.
        Identify each component ($C$, $I$, $G$, $NX$) clearly.

  \item Calculate GDP using the \textbf{income approach}.
        Identify all labour income and capital (profit) income.

  \item Calculate GDP using the \textbf{value-added approach}.
        Show the value added at each stage of production.

  \item Verify that all three methods give the same answer.
        If they do not, identify and correct any error in your work.
\end{enumerate}

\bigskip

%-----------------------------------------------------------------------
% PART II: INTERMEDIATE APPLICATIONS
%-----------------------------------------------------------------------

%-----------------------------------------------------------------------
\noindent\textbf{Problem 4} \quad \textit{Real GDP, nominal GDP, and the GDP deflator.}

\medskip
Use the Silverdale economy from Problem~3 as Year~1.
In Year~2, the same quantities of all goods are produced and sold, but:
\begin{itemize}[leftmargin=1.5em,itemsep=3pt]
  \item The price of timber rises to \$75 per cubic metre.
  \item The price of a finished plank bundle rises to \$11{,}500 (for the
        100-cubic-metre lot).
  \item The price of a chair sold domestically or to government rises
        to \$1{,}500; the export price rises to \$1{,}700 per chair.
\end{itemize}

\begin{enumerate}[label=(\alph*),leftmargin=2em,itemsep=4pt]
  \item What is nominal GDP in Year~2?
        (Hint: use the expenditure approach on final goods only.)
  \item What is \emph{real} GDP in Year~2, using Year~1 prices as the
        base year?  What is the real GDP growth rate between Year~1 and
        Year~2?
  \item Calculate the GDP deflator for Year~2 (with Year~1 as base year).
        What is the percentage change in the general price level?
  \item Explain briefly why real GDP growth and nominal GDP growth differ
        in this example.
\end{enumerate}

\bigskip

%-----------------------------------------------------------------------
\noindent\textbf{Problem 5} \quad \textit{Economic convergence.}

\medskip
In the year 2000, \textbf{Country Alder} has GDP per capita of
\$5{,}000 and grows at a constant rate of 4\% per year.
\textbf{Country Birch} has GDP per capita of \$24{,}000 and grows at
2\% per year.

\begin{enumerate}[label=(\alph*),leftmargin=2em,itemsep=4pt]
  \item Derive a formula for the year $t^*$ at which Country Alder
        catches up to Country Birch.  Solve for $t^*$ numerically.

  \item Suppose that, in the first 15 years, Country Alder instead
        grows at 7\% (due to institutional reforms), and then reverts
        to its original 4\% growth rate thereafter.  How many years after
        the year 2000 does it take Country Alder to catch up with
        Country Birch under this scenario?

  \item Give a brief economic interpretation of why the catch-up time
        changes.  What real-world mechanism might allow Country Alder to
        grow at 7\% in the early period?
\end{enumerate}

\bigskip

%-----------------------------------------------------------------------
\noindent\textbf{Problem 6} \quad \textit{From aggregate to per-capita variables.}

\medskip
Consider a closed economy with no government sector.  The aggregate
production function is:
\[
  Y_t = A\,K_t^{\,\alpha}\,L_t^{\,1-\alpha}, \qquad 0 < \alpha < 1.
\]
The population grows at rate $\bar{n}$: $L_{t+1} = (1+\bar{n})L_t$.
Households save a fixed fraction $\bar{s}$ of output each period:
$I_t = \bar{s}\,Y_t$.  Capital depreciates at rate $\delta$:
\[
  K_{t+1} = (1-\delta)K_t + I_t.
\]
Lowercase letters denote per-person quantities: $k_t = K_t/L_t$,
$y_t = Y_t/L_t$.

\begin{enumerate}[label=(\alph*),leftmargin=2em,itemsep=4pt]
  \item Show that $y_t = A k_t^{\,\alpha}$.
  \item Starting from the capital accumulation equation above, derive
        the per-capita capital accumulation equation:
        \[
          \Delta k_{t+1} = \bar{s}\,y_t - (\delta+\bar{n})\,k_t.
        \]
        Show every algebraic step clearly.
  \item Interpret each term on the right-hand side of the equation
        derived in part~(b) economically.  What does each term represent?
\end{enumerate}

\bigskip

%-----------------------------------------------------------------------
% PART III: MULTI-STEP ANALYTICAL PROBLEMS
%-----------------------------------------------------------------------


%-----------------------------------------------------------------------
\noindent\textbf{Problem 7} \quad \textit{Solow model: steady state.}

\medskip
Use the general setup from Problem~6 with $\alpha = 1/3$.

\begin{enumerate}[label=(\alph*),leftmargin=2em,itemsep=4pt]
  \item Find the steady-state capital per capita $k^*$ and output per
        capita $y^*$ in terms of the parameters $A$, $\bar{s}$, $\delta$,
        and $\bar{n}$.
  \item Suppose the savings rate $\bar{s}$ increases.  What happens to
        $k^*$ and $y^*$?  Use the formulas from part~(a) to explain,
        and give an economic interpretation.
  \item Suppose the population growth rate $\bar{n}$ increases.  What
        happens to $k^*$ and $y^*$?  Provide an economic explanation for
        the direction of the effect.
\end{enumerate}

\bigskip

%-----------------------------------------------------------------------
\noindent\textbf{Problem 8} \quad \textit{Solow model: numerical dynamics.}

\medskip
Use the per-capita capital accumulation equation from Problem~6(b)
with $\alpha = 1/3$ and the following parameter values:
\[
  A = 1,\quad \bar{s} = 0.25,\quad \delta = 0.06,\quad \bar{n} = 0.02,
  \quad k_1 = 6.
\]

\begin{enumerate}[label=(\alph*),leftmargin=2em,itemsep=4pt]
  \item Calculate $k^*$ and $y^*$ using your formulas from Problem~7(a).
  \item Calculate $\Delta k_2$, $\Delta k_3$, $\Delta k_4$, and
        $\Delta k_5$.  Show your work for at least the first two periods.
  \item Describe the pattern you observe in the sequence
        $\Delta k_2,\Delta k_3,\Delta k_4,\Delta k_5$.
        Explain \emph{why} this pattern arises.
  \item Is the economy above or below its steady state at period 1?
        How can you tell from your calculations?
\end{enumerate}

\bigskip
\newpage
%-----------------------------------------------------------------------
\noindent\textbf{Problem 9} \quad \textit{Solow diagram.}

\medskip
Using the parameter values from Problem~8, draw a fully labelled
\textbf{Solow diagram} with $k$ on the horizontal axis.

\begin{enumerate}[label=(\alph*),leftmargin=2em,itemsep=4pt]
  \item Label the two curves, their intercepts (if any), and the
        steady-state capital level $k^*$.
  \item Mark the initial position $k_1 = 6$ on the diagram.
  \item Draw an arrow or arrows indicating the direction of movement
        of $k_t$ from $k_1$ toward the steady state.
  \item What would happen to the diagram if the depreciation rate
        $\delta$ increased?  Describe which curve(s) shift and in
        what direction, and state the new qualitative effect on $k^*$.
\end{enumerate}

\bigskip

%-----------------------------------------------------------------------
% PART IV: SYNTHESIS
%-----------------------------------------------------------------------


%-----------------------------------------------------------------------
\noindent\textbf{Problem 10} \quad \textit{Comparative statics and policy analysis.}

\medskip
Return to the general Solow model from Problems~6 and~7.

\begin{enumerate}[label=(\alph*),leftmargin=2em,itemsep=4pt]
  \item A government implements a policy that raises the savings rate
        permanently from $\bar{s}$ to $\bar{s}' > \bar{s}$.
        Starting from the old steady state $k^*$, describe the
        \emph{transitional dynamics}: what happens to $\Delta k_{t+1}$
        immediately after the policy, and what happens as the economy
        converges to the new steady state ${k^*}'$?

  \item Using the expression for $y^*$ from Problem~7(a), show
        algebraically that a country with a \emph{higher} savings rate
        will have a \emph{higher} level of output per capita in the
        steady state, everything else equal.

  \item The Solow model predicts \emph{conditional convergence}.
        Explain what this means.  How does it differ from
        \emph{absolute convergence}?  Use the steady-state formula
        for $y^*$ in your answer.

  \item In the Solow model without technological progress, the
        long-run growth rate of output per capita is zero.
        Explain why.  Does this seem consistent with the real-world
        experience of developed economies?  Briefly discuss one
        extension to the model that can generate sustained per-capita
        growth.
\end{enumerate}

\bigskip\hrule\bigskip

\begin{center}
  \textit{End of Practice Problem Set 1}\\[4pt]
  \textit{Instructor solutions are available separately.}
\end{center}

\end{document}
